\documentclass[a4paper,11pt]{article}

\usepackage{revy}
\usepackage[utf8]{inputenc}
\usepackage[T1]{fontenc}
\usepackage[danish]{babel}
\usepackage{amsmath,amssymb,amsfonts,amsthm,textcomp}


\revyname{MatematikRevyen}
\revyyear{2026}
\version{0.1}
\eta{$3$ minutter}
\status{Færdig}

\title{Giver det mening?}
\author{Katrine '23, Lea '23 \& Carl '21}

\begin{document}
\maketitle

\begin{roles}
\role{S1}[Skuespiller] Forvirret studerende
\role{S2}[Skuespiller] Forvirret studerende
\role{TA}[Skuespiller] Øvelseslærer
\role{Dx}[Danser] Dansere
\end{roles}

\begin{sketch}

\scene{Vi er til øvelsestime i et kursus, der indeholder beviser. TA har lige færdiggjort et bevis på tavlen.}

\says{TA} Så har vi vist [Sætningen vælges senere. Afhænger af dansen]. Giver det mening? 

\scene{S1 og S2 forstår ikke}

\says{S1} Jeg tror ikke helt, jeg forstår stadig. 

\says{S2} Heller ikke mig, men jeg er også lidt mere en \textit{visual learner}.

\says{S1} Uh ja, det ville også hjælpe mig med en lidt mere visuel forklaring.

\says{TA} Okay.. Okay.. Jeg mener faktisk [forelæseren i kurset] plejede at gemmengå beviset mere visuelt tidligere år. Men så skal jeg lige bruge lidt hjælp. Kan I flytte bordere, mens jeg henter nogle af de andre TA's?

\scene{S1 og S2 er lidt forvirrede, men begynder at flytte bordene, mens TA forlader scenen. Musikken begynder}

\scene{Dansenummer starter: 4 Minutes af Justin Timberlake og Madonna}

\scene{Sang slut. Publikum klapper af den fede dans. S1 og S2 er meget begejstrede. TA high-fiver de andre TAs og siger tak for hjælpen. TA vender sig mod S1 og S2}

\says{TA} Giver det bedre mening nu?

\says{S1}[Går fra meget begejstret til forvirret] JA! ja! ja.. jahh.. well.. måske.. nej, nok egentlig ikke

\says{S2} Jaer.. måske hvis du lige dansede nogle eksempler, så vi kunne se det blive anvendt?

\end{sketch}
\end{document}
